% answers_short.tex —— 2003 年·春季北京卷·理科 · 简短答案速查
% 由 tools/make_short_answers.py 自 answers.tex 的 \gksolhead/\gkanswer 提取，
% 并与 2003/2003_数学真题_春季北京理_答案要点.md 逐题比对后生成，勿手改（改源 MD 后重跑）。

\begin{gknotice}
\textbf{使用说明}：本节只给最终答案便于快速核对，完整推导见解析版。答案为编写者独立推导，\textbf{非官方标准答案}。
\end{gknotice}

\gkpart{选择题}
\begin{center}\begin{tabular}{ccccccccccccc}
\toprule
题号 & 1 & 2 & 3 & 4 & 5 & 6 & 7 & 8 & 9 & 10 & 11 & 12 \\
\midrule
答案 & C & A & C & D & D & A & D & B & A & B & C & B \\
\bottomrule
\end{tabular}\end{center}

\gkpart{填空题}
\begin{itemize}[leftmargin=2.6em]
  \item[\textbf{13.}] $2/\sqrt3$。
  \item[\textbf{14.}] 140、85。
  \item[\textbf{15.}] $2\sqrt3$。
  \item[\textbf{16.}] $p/2$。
\end{itemize}

\gkpart{解答题}
\begin{itemize}[leftmargin=2.6em]
  \item[\textbf{17.}] $(2,3)$。
  \item[\textbf{18.}] 定义域 $\mathbb R\setminus\{\pi/4+k\pi/2:k\in\mathbb Z\}$；偶函数；值域 $[-1,1/2)\cup(1/2,2]$。
  \item[\textbf{19.}] (1) 证明见解析；(2) $16/\sqrt{17}$；(3) $16/3$
  \item[\textbf{20.}] (1) 88 辆；(2) 月租 4050 元；最大月收益 307050 元
  \item[\textbf{21.}] (1) 证明见解析；(2) $3\pi l^2/32$
  \item[\textbf{22.}] (1) $y^2=4x$；(2) 不能为正三角形；钝角时 $t\in(-\infty,-10/\sqrt3)\cup(2/(3\sqrt3),2\sqrt3)\cup(2\sqrt3,+\infty)$
\end{itemize}
